\section{Rendering --- current state}
\label{sec:rendering-current}

\begin{statusbox}
\textbf{Illumo ships world meshes, painter-correct 2D composition, cubemaps,
and explicit offscreen passes.} Live
frame path is acquire/enroll + tokens + submit via
\symbolref{Renderer}/\symbolref{IBackend}.
Drawables emit tokens through \texttt{AppendCommands}; GL draw submission is
confined to private \pathref{Illumo/Source/Rendering/OpenGL/} (plus window
bootstrap). Headless \pathref{Illumo/TestSupport/Include/Illumo/Testing/MockBackend.h}
powers \texttt{IllumoTests}. CanvasView
presentation is \textbf{a padded camera cache + bounded CPU RGB fade + one RGB
display texture/quad} with tiled multi-rect uploads (see canonical note in
\pathref{docs/architecture-consensus.md}).
No second real GPU API.
\end{statusbox}

\subsection{What actually draws frames today}
\begin{enumerate}
  \item Modules call \symbolref{DispatchDrawables} and add \symbolref{DrawableBase*}
        pointers to \symbolref{Scene} \textbf{layers}
        (\texttt{World} CanvasView, \texttt{UI} console/splash, \texttt{Debug} FPS).
  \item \symbolref{Illumo::render} pumps completed managed asset work, then
        drives \symbolref{Renderer}:
        \texttt{BeginFrame} $\rightarrow$ \texttt{RenderScene} $\rightarrow$
        \texttt{EndFrame} (swap).
  \item \symbolref{Renderer::RenderScene}:
        initial screen clear + viewport tokens; walk layers
        World $\rightarrow$ UI $\rightarrow$ Debug, using each layer's effective
        custom passes or ordinary screen rendering; each drawable
        \texttt{AppendCommands}; if a drawable returns false, legacy immediate
        \texttt{Draw()} still runs (hybrid fallback; production drawables return true).
        It captures window dimensions and the primary camera MVP once for that
        extraction; matching \symbolref{GameVisual} emitters reuse those values.
  \item \symbolref{IBackend::SubmitCommandQueue} $\rightarrow$
        \symbolref{GLDevice::ExecuteCommandQueue} resolves handles and issues GL.
\end{enumerate}

\subsection{Compiled scene extraction (D-E12/D-R25)}
The per-frame \symbolref{Scene} stays a non-owning drawable list.
\symbolref{SceneGraph} owns persistent SoA slots and intrusive hierarchy links,
with graph-local generational handles. Structural edits lazily compile a
preorder, parent indices, and subtree ranges. Local TRS edits set one dirty bit
and lower a watermark; one forward pass resolves affected world transforms.
World queries walk ancestors without consuming compiled dirtiness. Matrix
input decomposes to TRS; shear and projective matrices are lossy.

Nodes hold interned nonunique names, opaque 64-bit payloads, and ordered borrowed
attachments. Bounds cache by nonzero attachment revision; zero requests polling.
World boxes transform each attachment before union, and backward refit computes
subtree boxes. Unknown bounds fail open. A lazy binned-SAH query index has a
linear correctness fallback; warm queries still poll revisions in linear time.
The 4,096-entry change journal reports overflow explicitly.

\texttt{SceneGraphDrawable}, not the graph, derives from
\symbolref{DrawableBase}. It obtains one immutable snapshot per renderer frame
and consumes its matrices, bounds, and attachment pointers for collection,
depth, and color. Two graph-owned buffers are reused; ring reuse, attachment
invalidation, destruction, and graph teardown expire views. The adapter checks
the view before each callback. Owners invalidate before changing or retiring
borrowed content and preserve emitted token payloads until submission returns.
Mutation during extraction is rejected. All callbacks remain main-thread affine.

The shared shadow fit follows caster collection. Camera-missed subtrees skip
camera tests but retain potential off-camera casters, and depth relevance uses
snapshot bounds after fitting. This resolves the proposal's early-shadow-flag
assumption without changing the renderer's shadow policy. The temporary
\texttt{sceneSnapshotExtraction=0} switch selects guarded direct traversal.

Sibling order is part of the contract (D-E26):
\texttt{setParent(node, parent, insertBefore)} places a node immediately before
a current child of \texttt{parent}, or last when \texttt{insertBefore} is null,
and \texttt{getPreviousSibling} complements \texttt{getNextSibling}. The cycle
check and Reparented change record apply as for any reparent, and an unchanged
position records nothing. Scene documents store nodes in graph preorder, so
sibling order round-trips, and undo can restore a deleted node to its exact
position.

\subsection{SceneInstance (D-E18/D-E19)}
Scenes reach the renderer through \symbolref{SceneInstance} in
\texttt{Illumo::Content}, not through product code. It owns the document
state, one \symbolref{SceneGraph}, one \symbolref{MeshVisual} per renderable
component, and the asset references those visuals hold; stable node ids are
graph names. Programs add \texttt{drawable()} (and \texttt{skybox()} when the
environment names one) to their frame Scene and call \texttt{update()} once
per frame. Per-node edits (transform, flags, components, reparent with sibling
position, create, remove) validate first, change only the nodes they touch,
and invalidate graph snapshots before reconfiguring an attachment; recolors
and resizes that keep a node's component kinds reconfigure the visual in
place.
\begin{itemize}
  \item \textbf{Primitives:} solid shapes (rect, ellipse, triangle, cube,
        pyramid, sphere) share one immutable unit mesh per shape and instance,
        enrolled through \texttt{AssetManager::acquireMesh(MeshData)}; the
        primitive extent is the visual's model scale and its color the tint.
        Wire cube and wire sphere stay procedural.
  \item \textbf{Assets:} mesh, texture, atlas, and cubemap assets load through
        \symbolref{AssetManager} by normalized virtual path. A missing or
        failed asset draws a magenta placeholder and records a warning the
        program can read (\texttt{warnings()}).
  \item \textbf{Lighting:} the first enabled, visible \texttt{light} component
        in preorder supplies the one directional light; otherwise the
        environment sun does. Ambient color and shadow settings are pushed
        into every visual. 2D scenes are unlit.
  \item \textbf{Skybox:} the environment's cubemap asset is drawn by a
        \texttt{SkyboxVisual} that is a separate drawable, not a graph node.
        \pathref{SkyboxVisual.cpp} moved into the shared guest engine so
        every guest can use it.
  \item \textbf{Picking:} \texttt{pickRay} tests attachment local bounds from
        \texttt{SceneGraph::raycastCandidates} with rotation, scale and
        effective visibility, preferring the later node in preorder on ties
        (drawn on top in 2D). Optional pick proxies give empty, light and
        camera nodes a selectable box.
\end{itemize}
Loading 2,000 nodes takes about 3.8\,ms and 1,000 batched transform edits
about 0.8\,ms in Release (\texttt{Illumo.Content.Bench.Instantiate2000}).

IllEd's \texttt{EditorDocument} wraps one \texttt{SceneInstance} and is the
only mutation gateway; the editor adds no per-node attachments. IllMeshViewer
opens scenes through the same class, and IllumoGame's \texttt{render3dTest}
diagnostic is \pathref{IllumoGame/Scenes/render3d-test.ilsc} instantiated by
it. Journal overflow resynchronizes bindings without rebuilding nodes. Capture
fixtures use the adapter; CA storage and UI remain separate. See
\pathref{docs/scene-graph-v2-design.md},
\pathref{docs/scene-graph-v2-plan.md}, and
\pathref{docs/content-packages-and-scenes-design.md} for evidence and gated
follow-ups.

\subsection{Layers, styles, and drawables (D-R14, D-R18)}
The initial clear explicitly selects the screen framebuffer and full-window
viewport. Each ordinary layer reestablishes that target and viewport before
its drawables, including after an offscreen custom pass. The OpenGL binding
cache starts unknown on each submission, so the first screen bind is issued
even if the preceding submission left an offscreen framebuffer active.
Pass color and depth clear flags are independent. Color-only passes emit a
color-buffer clear; depth clears carry the requested value. Combined pass
clears emit both commands. Clear tokens default to depth one, so a later
default clear does not inherit a custom value from a previous pass.

Scene stores host default passes separately from application overrides.
Nonempty application overrides take precedence, whether installed
in module Start, Update, or DispatchDrawables. The host updates only defaults
through \texttt{SetDefaultLayerPasses}, so motion-blur toggles and parameter
changes cannot erase application passes. An empty override,
\texttt{ClearLayerPasses}, or \texttt{ResetDefaultPasses} restores the current
host fallback; without one, rendering uses the ordinary screen draw.
\texttt{passesIn} and \texttt{hasCustomPasses} describe the effective sequence.
Required-module transitions reset all application overrides before destroying
retired modules and after rejected startup, detaching borrowed callbacks while
preserving defaults. Optional modules reinstall their overrides after a transition.
\begin{itemize}
  \item \textbf{Layer} --- ordered composition bucket, rendered to the screen
        by default or through its effective custom pass sequence.
  \item \textbf{RenderStyle} --- Renderer-owned generational registry: shader handle +
        \texttt{PipelineState} defaults (\texttt{Canvas}, \texttt{UiText},
        \texttt{Console}, \texttt{Shape}, \texttt{Sprite}). Canonical Shape/Sprite
        programs consume \texttt{uMVP} only (\symbolref{WorldLook}). Overlay
        chrome supplies a Y-down screen ortho; world objects supply camera
        view-projection times node world. GPU programs still live in backend
        registries.
  \item \symbolref{DrawableBase} --- type-erased list entry; virtual
        \texttt{Draw} + virtual \texttt{AppendCommands} (default returns false).
        \textbf{No GL handle ownership} on the base (D-R6). Content handles
        only; call \texttt{bindStyle} then emit content tokens.
  \item \symbolref{Drawable<Derived>} --- CRTP: \texttt{Draw} $\rightarrow$
        \texttt{DrawImpl} for the rare immediate path. Production emitters
        (CanvasView, CommandLine, GLString/SplashText) use
        \texttt{AppendCommands} only; \texttt{DrawImpl} is empty/no-op.
  \item \symbolref{GameVisual} (D-R15) hosts composed
        shape/sprite/text values. Parent and local \texttt{Transform2D},
        normalized pivots, atlas regions/flips, and integer draw order build one
        stable cross-type stream. Only adjacent compatible items batch, so
        transparent painter order is never changed by a global texture sort.
        Pixel-space rebuilds conservatively reject quads outside the logical
        viewport. An optional top-left logical-pixel clip rejects fully
        excluded quads before upload and uses nested, intersected scissor
        tokens for partial content; popping the clip restores any outer
        scissor (D-R23). World-space behavior is unchanged.
        Dynamic buffers reserve 1,024 quads, double as needed, and stop at a
        configurable 65,536 default ceiling.
  \item \symbolref{MeshVisual} (D-R21/D-R24) is the world mesh host and
        \symbolref{ISceneRenderAttachment}: colored lines/triangles, textured
        quads (sprites), optional camera-facing billboards, optional
        directional lighting with participation in a scene-level depth-only
        shadow pass, and optional object motion blur from the previous clip
        pose. Bounds cover every geometry source immediately after mutation;
        camera-facing sprites use a view-independent enclosing volume. It clones
        Shape/Sprite styles with depth-tested pipelines and emits
        \texttt{uMVP}. Lit meshes also emit \texttt{WorldLook} light, shadow,
        and motion-blur uniforms from CPU-side drawable state; products
        persist those values through EnvVars rather than having
        \symbolref{MeshVisual} read configuration.
        One host can contain multiple visual items. Use child scene nodes for
        independent transforms and subtree state; each node may borrow an ordered
        list of attachments. Overlay HUD stays on \symbolref{GameVisual}. Dynamic
        procedural meshes upload dirty vertices through \texttt{UpdateBuffer};
        capacity growth replaces the backing allocation. Immutable imported
        meshes instead borrow one reference-counted \texttt{MeshHandle} from
        \symbolref{AssetManager}; multiple visuals reuse the upload while
        retaining independent transform, tint, lighting, and motion state.
        The Renderer retains caster descriptors, chooses light settings from
        the first camera-visible caster, and filters them against a conservative
        camera-frustum extrusion toward the light. Invalid reconstruction falls
        back to all casters. A monotonically increasing frame serial resets
        previous-MVP history when a culled visual re-enters.
\end{itemize}

\subsection{Managed assets (D-R17/D-R24)}
Shared fonts cache one atlas handle per live renderer ownership identity,
not per raw renderer address. The identity is weak in the font cache; expired
entries are pruned and explicitly retired texture handles are reenrolled.
Alternating live renderers reuses their separate enrollments. Renderer teardown
invalidates the identity before backend teardown; GPU ownership stays with the
backend, and a font never dereferences a retired renderer.

\symbolref{AssetManager} caches canonical file paths plus loading options and
returns stable typed handles. Duplicate acquisition increments a reference
count; the final release destroys the backend resource. Initial pending or
failed 2D texture/shader requests retain a checkerboard or fallback shader that follows
the documented custom-2D binding contract. One owned worker performs file reads
and stb decoding; \texttt{pump} applies GPU
replacement on the render thread before frame submission. Failed reloads keep
the last good object and record an error. Debug builds poll timestamps every
500\,ms and coalesce requests; explicit reload is available in all builds.
Shutdown cancels obsolete results and joins the worker before backend teardown.

Static mesh acquisition is synchronous and main-thread affine. The
\symbolref{MeshLoader} converts a file to CPU \texttt{MeshData};
\symbolref{AssetManager} canonicalizes the file path plus geometry-affecting
options, enrolls one immutable backend mesh, and returns its typed handle.
Duplicate acquisition increments a reference count and final release destroys
the backend resource. In-memory mesh data can be enrolled once and explicitly
retained for multiple consumers. Managed mesh metadata supplies the immutable
vertex/index counts and bounds needed by drawables. Mesh hot reload, materials,
and automatic instance aggregation remain outside this milestone.

The byte source is an \texttt{IAssetSource}: the file system natively,
\texttt{VfsAssetSource} over the host file tree in tools and native tests, and
\texttt{GuestVfsAssets} in guests, where canonical keys are normalized
absolute virtual paths (D-E24). \texttt{AssetManager::assetSource()} exposes
the source so a product can read a scene through it. MTL libraries beside an
OBJ resolve through the same source: \texttt{MeshLoader::materialLibraryNames}
is the single \texttt{mtllib} scanner, \texttt{AssetManager::acquireMesh}
reads each named library relative to the OBJ, and the text reaches the loader
as \texttt{MeshLoadOptions::materialText}. Materials are parsed but still not
bound.

Cubemaps retain their source kind and all canonical face paths (or the cross
path). Initial acquisition remains synchronous; reload decoding uses the same
CPU job path and publishes all six RGBA faces through cubemap replacement.
Initial success has revision one. Explicit path reload and timestamp polling
observe every face. Decode timestamps are captured before reading, so a later
edit remains observable on the next poll. Missing/mismatched faces or rejected
replacement keep the prior ready resource and revision. Cross extraction and
orientation are shared by initial load and reload; texture kinds cannot change
during replacement.

\subsection{CanvasView presentation (bounded RGB fade + dirty rects)}
\label{sec:canvas-rgb-fade}
\begin{itemize}
  \item Domain: \symbolref{SparseCellGrid} stores signed-coordinate
        16$\times$16 chunks and changes its revision only for a cell-content
        change. The view visits only chunks intersecting its camera source
        bounds.
  \item View: visible diagnostics are separate from a globally aligned cache
        padded by two 16-cell chunks on every side. Motion inside the cache
        changes only the camera MVP. CPU palette via \texttt{evalCell}
        $\rightarrow$ \texttt{targetRgb};
        \texttt{displayRgb} eases toward targets, newly revealed cells snap,
        and packed RGB is staged for upload. Unchanged grid, palette, and
        camera bounds skip resampling entirely.
  \item Zoom: normal views retain one nearest-filtered texel per 16$\times$16
        world cell. Far views use a stable integer density LOD bounded to about
        four screen pixels per texel; it coarsens immediately to fit and refines
        only when the next level fits within 80\% of the budget. Changed chunks
        map to deduplicated exact or overview cache bins.
  \item GPU: one reusable nearest-filtered RGB display texture and one
        world-space quad aligned to the same 16$\times$16 cell bounds as the
        editor cursor. Dirty 16$\times$16-texel tiles merge into at most eight
        \texttt{UpdateTexture} rectangles when their area is at most half the
        enclosing AABB; otherwise the AABB is submitted. Small dimension
        increases reserve 50\% texture headroom. Re-enrollment destroys the
        prior GL texture, PBOs, and fences under the same handle, and view
        destruction explicitly releases the texture. Upload-run and output-rect
        scratch vectors are retained across dirty frames.
  \item Shader/style: \symbolref{CanvasView} composes a world-space
        \symbolref{GameVisual} sprite and therefore uses the Renderer-owned
        inline \texttt{Sprite} style with sampler \texttt{uTexture}.
        The file-backed \texttt{Canvas} style and
        \pathref{Shader/canvas\_*.glsl} remain enrolled but have no production
        CanvasView binding.
  \item History: D-P4 briefly used GPU R8 + $256\times 1$ palette (snap colors).
        Fade was restored; dirty-rect PBO + bind-state tracker from that work
        remain. Canonical write-up:
        \pathref{docs/architecture-consensus.md} \S5.6.
\end{itemize}

\subsection{Backend surface}
\symbolref{IBackend}: lifecycle and command-queue push/submit/clear. Resource
creation returns non-convertible slot+generation \texttt{MeshHandle},
\texttt{ShaderHandle}, and \texttt{TextureHandle} values. Validated
replace/destroy/query operations and command resolution reject stale handles
with a warning and safe no-op.

\symbolref{CommandQueue}: vector storage reserves 2,048 tokens, grows to a
configurable 65,536 default safety ceiling, and reports high-water and rejected
counts. Rejection logs once per reset interval. On the supported 64-bit Windows
build each token is 72 bytes. Matrix-uniform values are copied into retained
renderer-owned chunks of 128 matrices and the compact token borrows that stable
value until synchronous submission returns; the chunks reset after submission
and are reused by later frames. Matrix side storage uses a matching default
65,536-value safety ceiling so rejected token storms cannot grow it without
bound.

\symbolref{RenderCommand}: tagged union (\texttt{CommandType} + POD payload)
including pipeline, bind, uniforms, \texttt{UpdateTexture}, \texttt{UpdateBuffer},
clear, draw, and explicit enabled/disabled scissor state (D-R1).

Each \symbolref{GLShaderProgram} owns its uniform-location cache. Cache hits use
the token's name directly rather than building and hashing a combined
program-ID/name string, and shader replacement or destruction retires cached
locations with the owning program.

Renderer maintains a per-frame scissor stack for \texttt{pushClipRect}/
\texttt{popClipRect}. Nested clips intersect in physical viewport coordinates
and restore the prior state; low-level \texttt{pushScissor} remains the explicit
state-token path (D-R23).

\symbolref{GLTexture::UpdateSubImage}: rectangles through 64\,KiB use direct
\texttt{glTexSubImage2D}. Larger uploads select a signaled slot from a
three-PBO/\texttt{GLsync} ring without waiting; all-busy and map-failure cases
fall back to the direct path. Replacement/destruction deletes every fence and
buffer and retains the OpenGL 3.3 baseline.

Uploads reject invalid bounds, channel counts, row strides, and byte-size
overflow before any staging or GL call. Uploads bind on texture unit zero,
matching the device cache, and restore the caller's active unit and unpack
state. Cubemap file decoding explicitly produces RGBA for every face and cross;
raw two-channel cubemaps are rejected instead of reinterpreted as RGBA. OBJ
conversion validates all attribute references before publishing mesh data.

\symbolref{GLDevice}: bind-state tracker (skip redundant program/VAO/texture/
viewport within a submit); uniform location cache; blend-func-on-enable (D-R5).

Presentation defaults to monitor synchronization through
\texttt{glfwSwapInterval(1)}. The persisted \texttt{vsync} value selects that
frame-paced mode or an explicit uncapped profiling mode and may change live at
a swap boundary. Debug labels swap-completion cadence as \texttt{Paced FPS} and
main-loop work as \texttt{Submit FPS}; uncapped mode reports paced FPS as off
instead of treating thousands of CPU submissions as displayed frames (D-P32).

\subsection{Primitive-composed UI (D-R20)}
\begin{itemize}
  \item \symbolref{CommandLine}: one batched \texttt{DrawIndexed} when open.
        \texttt{UpdateBuffer} runs when history, input, scroll, layout, caret
        phase, or a quantized accent pulse changes. Wrap metrics are cached
        until contents or panel width change; only visible history lines are
        tessellated. Its surface, inset history, input row,
        selection, caret, scrollbar, and floating resize affordance compose
        \symbolref{GameVisual} fills, outlines, lines, and text.
  \item \symbolref{GLString}: geometry cache --- skip VBO re-upload if
        content/style unchanged (FPS $\sim$1\,Hz rebuilds). Optional panel
        chrome adds a shadow, surface, border, and accent in the same batch.
  \item \symbolref{SplashText}: fading EDIT/NORMAL label built on decorated
        GLString chrome; the two modes use distinct shared theme accents.
  \item \texttt{UiTheme}: shared value-only colors and compact panel metrics;
        it is not a retained widget tree or layout engine.
  \item \symbolref{GameVisual} gradient quads (D-R26): convex, fan-ordered
        quads with one color per vertex through the ordinary shape vertex
        format and shader. GuiKit composes soft chrome from them (glass
        panels, gradient surfaces, soft shadows, radial glows, vignettes,
        sheens, keycaps); fades keep their color at zero alpha so
        straight-alpha blending never darkens an edge.
  \item \symbolref{DebugModule}: \texttt{renderer\_demo}, \texttt{assets}, and
        \texttt{asset\_reload} prove managed atlas/shader loading, atlas
        animation, transforms, tint, alpha order, async load, state inspection,
        and reload in Debug builds only.
\end{itemize}

\subsection{Tests (headless)}
\texttt{IllumoTests} owns persistent scene hierarchy plus reusable
rendering/services cases, while
\texttt{IllumoGameTests} owns CanvasView, ruleset, CellGameModule, simulation,
and persistence cases. Editor and mesh-viewer cases have their own runners.
These are shared binaries for compile efficiency, but
CTest invokes one exact logical case per isolated process (D-T1). No GPU
context is required. Combined Clang/LLVM \texttt{IllumoCoverage} enforces 85\%
line coverage over headless-testable first-party production code linked into
all registered workspace runners.

\subsection{Remaining pain points (non-blocking)}
\begin{enumerate}
  \item Hybrid immediate \texttt{Draw()} remains for \emph{test stubs only};
        production drawables are pure-token (D-R10).
  \item String-named uniforms are GL-shaped debt if a second real GPU API appears
        (D-R9).
  \item \pathref{RenderWindow} still uses GLEW for context create / viewport
        (window bootstrap, not frame submission).
  \item Windows is the supported live path. Linux window bootstrap no longer
        issues GL calls before GLEW; native GUI smoke is still required.
        macOS is not targeted and has no retained scaffold.
  \item Font atlases/UTF-8 layout, tilemaps/particles, mesh materials and formats
        beyond OBJ, and broader offscreen compositing are follow-on features,
        not hidden scaffolding in the current renderer.
\end{enumerate}


\subsection{Standalone frame capture (2026-09-11)}
\texttt{FrameCapture} creates a hidden, unpaced OpenGL context for one synchronous
producer callback. No application host, input, simulation, audio or persistent
configuration starts. The producer owns its content through normal token
submission and destroys renderer-bound objects before returning. The renderer
reads top-down RGBA8 pixels before swap and tears down before the context.
Strict capture rejects immediate fallback, invalid resources and queue overflow.
Existing applications retain normal scheduling and visible-window defaults.
The former \texttt{IllumoCapture} fixture CLI is removed (D-E14). Command-line
capture is now \texttt{IllumoRuntime --capture}: a running WASM application
renders to \texttt{--capture-frame} (default 60), a
\texttt{Renderer::setBeforePresent} hook reads back the presented backbuffer,
\texttt{FrameCapture::savePng} publishes a new PNG, and one JSON result line
reports success, size and error. \texttt{--capture-script} first runs
bench-script steps (console lines, key presses, waits) so a capture can reach a
later screen, and the saved frame is forced opaque to match the presented
window. \texttt{FrameCapture} itself remains the
engine API, checked on a real GPU by \texttt{IllumoCaptureGpuTests}.
See \texttt{docs/frame-capture.md} for inputs, publication and limitations.

Current 3D facilities include CPU OBJ import, directional lighting,
per-drawable self-shadowing and post-processing, and scene-level environment
lighting, sun and skybox resolved by \texttt{SceneInstance} from a scene's
\texttt{settings.environment}. These do not yet implement registered game
objects or physics integration.

Renderer-bound GameVisual resources retain one renderer lifetime. Foreign use,
truncated geometry, and queue rejection report incomplete frames and suppress
presentation. Rejected retained uploads retry on a subsequent frame. Backend
queue totals survive pass resets; native GPU names live only in private GL types.
Mesh and texture enrollment validates actual storage before publishing a handle;
failed replacements preserve old resources and oversized vertex writes fail early.
